\section{Model Evaluation} \label{sec:model-evaluation}
After training, we evaluated model performance on the test set clips to ensure that seizure detection was substantially better than chance. 
Only under this condition do the \glspl{fp} provide meaningful information about the model’s seizure propensity; \glspl{fp} from a near‑random classifier would not.


The models were employed to assign scores in the range [0, 1] to all existing segments, representing the model-derived probability that a given segment is preictal.  
Subsequently, the segment-level scores were averaged within each clip to obtain a corresponding clip-level score.  
These clip-level scores, rather than the raw segment-level scores, were used in subsequent evaluations.
Two types of metrics were used to evaluate the performance: clip-based metrics and \glspl{ebm}.

Here, clip-based metrics refer to the standard measures used to evaluate models in machine learning on samples (clips), which use the confusion matrix of \glspl{tp}, \glspl{fp}, \glspl{fn}, and \glspl{tn} to compute metrics, such as precision, sensitivity, accuracy, and the \gls{roc auc}.
The Hanley McNeil test \parencite{hanley_method_1983} was performed to test for statistically significant improvement of the \gls{roc auc} compared to a random classifier.

On the other hand, \glspl{ebm} are specific to time-series data and seizure prediction and quantify more clinically relevant properties.
Here, the (relative) \gls{tifw} \parencite{mormann_seizure_2007, eberlein_machine_2025} and the (relative) number of seizures predicted were used.

\subsection{Time in False Warning} \label{subsec:tifw}
Relative \gls{tifw} is the \gls{eb} equivalent to the \gls{fpr}, i.e., the inverse of specificity.
That is to say it is minimized when non-preictal (false) samples are correctly classified as such.
However, when a prediction is issued from a $\sim$\SI{10}{\minute} clip, it is not issued for the timespan of that clip, but for that clip's \gls{sph}, i.e., the timespan $\sim$\num{5}--\SI{65}{\minute} after the clip's end.
This means that summing the duration of \gls{fp} clips is not sufficient to calculate the \gls{tifw}.
The clip's \glspl{sph} and overlaps thereof must be taken into account.
When a system unnecessarily alarms the patient, it causes them distress and inconvenience, which is why it is essential to minimize false alarms.

To calculate \gls{tifw}, only valid, non-preictal clips were selected.
For all such clips that issued a warning (\glspl{fp}), the \glspl{sph} were selected.
Overlapping \glspl{sph} were then merged, so that shared time was not counted multiple times.
The total duration of these non-overlapping false-warning intervals was the absolute \gls{tifw}.
The relative \gls{tifw} was obtained by normalizing absolute \gls{tifw} by the maximum possible warning time, which was the sum of the non-overlapping \gls{sph} intervals from all non-preictal clips.
This yielded a metric bounded between 0 and 1, enabling comparable interpretation across patients and consistent with the guidelines in \textcite{mormann_seizure_2007}.

\subsection{Event-Based Sensitivity} \label{subsec:eb-sensitivity}
It is also essential to quantify the relative number of seizures successfully predicted, which is also referred to as \gls{eb} sensitivity.
For preictal clips, a positive prediction is considered a true warning, and a seizure is counted as predicted if it falls within the \gls{sph} of at least one positively predicted preictal clip.
This yields the number of predicted seizures, and \gls{eb} sensitivity is then computed as the ratio of predicted to total seizures.
This metric therefore reflects clinically relevant event detection performance at seizure level rather than clip level.

\subsection{Threshold Optimization} \label{subsec:threshold-optimization}
Whether a clip issued a warning depended on whether its score was greater than or equal to a classification threshold.
To optimize it, metrics were calculated for all possible thresholds, i.e., all unique model scores.
To take both the \gls{eb} specificity (1 - rel. \gls{tifw}) and the \gls{eb} sensitivity into account, the harmonic mean between them was calculated:
\begin{equation} \label{eq:harmonic-mean}
    H(a,b)=\frac{2ab}{a+b}
\end{equation}
\begin{equation} \label{eq:eb-score}
    S_{\mathrm{EB}} = H\!\left(\mathrm{Sensitivity}_{\mathrm{EB}},\, 1-\mathrm{rel.\ TIFW}\right)
\end{equation}
The threshold that maximized this \gls{eb} score was set as the best threshold and used to determine binary clip predictions in further steps.

The metrics were calculated independently on the train and test set.
In causal, (pseudo-)prospective study designs, classification thresholds should be determined exclusively using the training set. 
However, as this work prioritizes the cycle-based evaluation over a realistic measure of performance, the threshold was individually optimized for the train and test set, which boosted performance on the test set.
